\section{When a sparse edit budget cannot repair the anchor}
\label{supp:sparse-anchor-risk}
We compare the mean and minimum of nine leave-one-fitting-source-out cost
heads, for three seeds. Both policies choose between the same two recipients
per query, with the same peak edit relative to OLS4 and with directions fixed
by the full-fit head. At 1,613 matched edits per seed, minimum scoring improves
observed utility by \SparseKnown{} points (95\% source-bootstrap interval
\SparseKnownCI{}) over mean scoring, but its full-grid lower contrast is
\SparseLower{} points. Neither passes the original 2\% rule in any of the
102 repeated recording views. The minimum of these dependent fits is a
sensitivity score, not a calibrated uncertainty bound.

We ask whether \emph{any} one-edit-per-query policy starting at OLS4 could
pass the CV-relative rule on these recordings. For all or easy targets,
let $h_i=(e_i^{\mathrm{OLS4}}-e_i^{\mathrm{CV}})_+$ on known rows in that
subset and zero elsewhere. With arbitrary edits, the policy can eliminate
at most $\max_{i\in q} h_i$ known harm in query $q$; adding further harm is
unnecessary for this optimistic calculation. The actual risk denominator
cannot exceed the known CV error mass $M$ over the whole subset. For $M>0$,
an optimistic lower bound on the original positive-harm ratio is
\[
  \frac{\sum_i h_i-\sum_q\max_{i\in q} h_i}{M}.
\]
Unknown harm is set to zero, and all and easy can erase different rows. The
bound exceeds 2\% on \SparseAll{}/\SparseN{} recordings for all targets,
\SparseEasy{}/\SparseN{} for easy, and \SparseUnion{}/\SparseN{} for their
union. This is a necessary-condition check for an OLS4 anchor at this budget,
not a change to the safety rule. It shows why a strong accuracy comparator
need not be a suitable deployment anchor. No independent outcomes were used.
